\section{Objective}
\label{sec:6-1}

Determine what the transaction data reveal about sales, customers, revenue composition, repeat purchasing, behavior, and the feasibility of churn prediction.

\textbf{Note.} \textbf{Iterative link with Chapter~\ref{ch:data-preparation}.} Sections~\ref{sec:6-13}--\ref{sec:6-14} need churn labels and snapshots, which are formally built in Chapter~\ref{ch:data-preparation}. This chapter uses \textbf{provisional} labels and snapshots computed for candidate values of $H$, $R_{\max}$ and $C$. Final values are fixed in Chapter~\ref{ch:data-preparation}, following the iterative nature of the Data Science methodology.

\section{Dataset Overview}
\label{sec:6-2}

Transaction count, unique customers, unique products, date range, countries, total quantity and revenue.

\section{Data Quality}
\label{sec:6-3}

Missing values (especially Customer ID), duplicate records, invalid Quantity/Price values, cancellation records, non-product stock codes and other quality issues.

\section{Monthly Sales Performance}
\label{sec:6-4}

Monthly customers, orders, quantity and revenue to understand sales dynamics and seasonality.

\section{New vs Existing Customers}
\label{sec:6-5}

Classify customers by whether the period contains their first valid purchase or a subsequent one. Compare counts and revenue contribution.

\section{Revenue from New and Existing Customers}
\label{sec:6-6}

New Customer Revenue, Existing Customer Revenue, Total Revenue and monthly shares. This directly supports the business case for retaining existing customers.

\section{Repeat Customer Analysis}
\label{sec:6-7}

One-time versus repeat customers: counts, shares, orders and revenue contribution.

\section{Monthly Repurchase Analysis}
\label{sec:6-8}

Eligible existing customers, repeat buyers and repurchase rate by month.

\section{Inter-Purchase Interval}
\label{sec:6-9}

Time between consecutive purchase occasions: mean, median, percentiles and distribution. Results inform \textbf{$H$} and \textbf{$R_{\max}$}.

\section{Customer Monetary and Behavioral Analysis}
\label{sec:6-10}

Customer revenue, order frequency, quantity, average order value and behavioral heterogeneity (including the influence of large wholesale customers).

\section{Customer Lifecycle}
\label{sec:6-11}

Progression from first purchase to repeat purchase, active behavior, increasing recency, at-risk behavior and potential churn.

\section{RFM Exploration}
\label{sec:6-12}

Recency, Frequency and Monetary distributions and their relationship with subsequent purchasing.

\section{Churn Exploration}
\label{sec:6-13}

Using provisional labels: churn rate over time and across recency, frequency, monetary value and other characteristics. \textbf{Report the overall churn rate early} -- with many one-time buyers, churn may be the \emph{majority} class, which changes how PR-AUC and Precision should be read.

\section{Snapshot Population and Design Feasibility}
\label{sec:6-14}

\begin{enumerate}
    \item Compare daily, weekly-all and weekly episode-based snapshots: number of rows, rows per customer, share of near-duplicate rows, and class distribution. This quantifies why daily snapshots are rejected and how much the episode rule reduces training redundancy.
    \item Show the effect of the activity condition ($R_{\max}$) on population size and churn rate.
    \item Report the \textbf{weekly Scoring Population size} over time, which is the pool from which marketing selects Top-K each week, and contrast it with the much smaller number of customers that would be scored if the episode rule were wrongly applied to scoring.
    \item \textbf{Feasibility check:} for each candidate $H$, count labelable weekly scoring dates after the warm-up and after the embargoes in Section~\ref{sec:7-9}. The test period should contain enough scoring dates (e.g.\ $\geq 8$) for stable Top-K estimates.
\end{enumerate}

Use the evidence to justify weekly batch scoring, the training-episode cooldown, $R_{\max}$ and the final $H$.

\section{Key Findings}
\label{sec:6-15}

Business-oriented findings on revenue composition, repeat purchasing, customer heterogeneity, purchase intervals, churn behavior, data limitations and the chosen snapshot/prediction-window design.
